\section{From algebraic maneuvers to maps and relations}
\label{f:begin}

When we write $a=b+c$ and then $b=a-c$, something more organized than
symbol manipulation is happening. We have applied an invertible map. When
we then write $d(a-c)=da-dc$, we have used a map that preserves an
operation. These two ideas---invertibility and preservation---will remain
visible as the objects become more general.

\subsection{An elementary maneuver is a diagram}

Let $(V,+,0)$ be an abelian group. For $c\in V$, translation is the map of
sets $\tau_c:V\to V$, $x\mapsto x+c$. Associativity, the identity law,
and the inverse law give
\[
\tau_{-c}\tau_c(x)=(x+c)+(-c)=x+(c-c)=x.
\]
The same calculation in the other order proves
$\tau_{-c}=\tau_c^{-1}$. Thus $a=b+c$ means $a=\tau_c(b)$, and solving
for $b$ means applying the inverse. Translation need not preserve the
zero element, so it is generally an isomorphism of sets, not an
automorphism of abelian groups.

Associativity itself says that two composites are the same map:
\[
\begin{tikzcd}[column sep=large]
V\times V\times V \arrow[r,"+\,\times\,\id"]
  \arrow[d,"\id\,\times\,+"'] & V\times V \arrow[d,"+"]\\
V\times V \arrow[r,"+"'] & V.
\end{tikzcd}
\]
Here the usual canonical identification of the two parenthesizations of
the triple product is understood. The diagram says that making one
intermediate sum rather than the other does not change the result.

If $W$ is another abelian group and $h:V\to W$ preserves addition,
then $h(0)=0$: the equality
$h(0)=h(0)+h(0)$ can be translated by $-h(0)$. Also
$h(-c)=-h(c)$, because $h(c)+h(-c)=h(0)=0$. Consequently
\[
h(a-c)=h(a)-h(c),\qquad h\tau_c=\tau_{h(c)}h.
\]
For a module, multiplication by a scalar is an additive homomorphism,
so distributivity is exactly this preservation statement. In a group,
the same discussion uses multiplication and left or right translations,
with the order retained. Its multiplication is a map of underlying
sets; it is not generally a homomorphism from the product group.

There is also a structural explanation of the triangular substitutions
used in basis exchange. Given an additive homomorphism $u:V\to W$, the
\emph{shear}
\[
S_u:V\times W\longrightarrow V\times W,
\qquad (v,w)\longmapsto(v,w+u(v))
\]
has inverse $S_{-u}$. This follows by composing the maps, without counting
any basis. Over a field, if $t_1\ne0$, the map
\[
(x_1,\ldots,x_n)\longmapsto
(t_1x_1,t_2x_1+x_2,\ldots,t_nx_1+x_n)
\]
is a scaling in the first coordinate followed by the shear $S_u$ on
$\kk\times\kk^{n-1}$, where
\[
u:\kk\longrightarrow\kk^{n-1},\qquad
u(z)=(t_2t_1^{-1}z,\ldots,t_nt_1^{-1}z).
\]
Both maps are invertible. The inverse of the composite sends $y_1$ to
$t_1^{-1}y_1$ and $y_i$ to
$y_i-t_i t_1^{-1}y_1$ for $i>1$. This establishes the elementary maneuver
without presupposing invariance of dimension.

The aim of the following constructions is similar: replace a list of
rearrangements by maps whose behavior is controlled once, and then reuse
that control.

\subsection{What is retained by a universal construction?}

A \emph{category} consists of objects, sets of maps between objects,
identity maps, and associative composition. A diagram commutes when
every indicated pair of composites with the same endpoints agrees.
A map $m$ is a \emph{monomorphism} if $mu=mv$ implies $u=v$; a map $e$
is an \emph{epimorphism} if $ue=ve$ implies $u=v$. These are cancellation
conditions on maps, not definitions in terms of elements.

A \emph{subobject} of $X$ is a monomorphism into $X$, with two such
monomorphisms identified if an isomorphism over $X$ relates them. We
write $A\leq B\leq X$ when the map representing $A$ factors through the
one representing $B$. Products are characterized by their projections:
a map to $X\times Y$ is uniquely a pair of maps to $X$ and $Y$.

\begin{definition}[Pullback]
For $f:X\to Z$ and $g:Y\to Z$, a pullback is an object $P$ with maps
$a:P\to X$, $b:P\to Y$ such that $fa=gb$, universal with that property:
every pair $u:T\to X$, $v:T\to Y$ with $fu=gv$ has a unique map
$t:T\to P$ satisfying $at=u$ and $bt=v$.
\[
\begin{tikzcd}
P \arrow[r,"b"] \arrow[d,"a"'] & Y \arrow[d,"g"]\\
X \arrow[r,"f"'] & Z.
\end{tikzcd}
\]
We write $P=X\times_ZY$.
\end{definition}

In sets this is the set of pairs with the same image in $Z$. In groups
and modules it has the same description with coordinatewise operations.
The universal property retains the equation $fu=gv$, without selecting
coordinates. A terminal object receives exactly one map from each object.
An initial object has exactly one map to each object.
An equalizer of $u,v:X\to Y$ is universal among maps $a:T\to X$ with
$ua=va$. A finite-limit category has a terminal object, finite products,
and equalizers, hence all the pullbacks used here: $X\times_ZY$ is the
equalizer of $f\pi_X,g\pi_Y:X\times Y\to Z$. The
\emph{kernel pair} of $f:X\to Y$ is
$X\times_YX\rightrightarrows X$: it records exactly which inputs $f$
identifies.

A \emph{coequalizer} of $u,v:K\rightrightarrows X$ is a map $q:X\to Q$
with $qu=qv$ such that every map $h:X\to Z$ satisfying $hu=hv$ factors
uniquely as $h=\bar hq$. A \emph{regular epimorphism} is a coequalizer
of some pair. Every regular epimorphism is an epimorphism, by the
uniqueness part of this property. We draw it as $X\twoheadrightarrow Q$.

\begin{lemma}[The elementary descent rule]\label{f:descent}
In a category with pullbacks, a regular epimorphism $q:X\twoheadrightarrow Q$
is the coequalizer of its kernel pair. It has the following lifting
property against any monomorphism $m:A\hookrightarrow B$: in a
commutative square
\[
\begin{tikzcd}
X \arrow[r,"a"] \arrow[d,two heads,"q"'] & A \arrow[d,hook,"m"]\\
Q \arrow[r,"b"'] \arrow[ur,dashed,"d"] & B,
\end{tikzcd}
\]
there is a unique $d$ with $dq=a$ and $md=b$.
\end{lemma}
\begin{proof}
Choose a pair $u,v$ of which $q$ is a coequalizer. Since $qu=qv$, this
pair factors through the kernel pair of $q$. A map equal on that kernel
pair is therefore equal on $u,v$, and descends uniquely through $q$.
This proves the first claim.

For the square, $ma u=bq u=bq v=ma v$, and monicity of $m$ gives
$a u=a v$. Hence $a=dq$ for a unique $d$. Now $mdq=ma=bq$, so
$md=b$ because $q$ is epic. The equation $dq=a$ also proves uniqueness.
\end{proof}

In particular, a regular epimorphism that is also a monomorphism is an
isomorphism: apply the lifting property to its square with itself and
identity maps. This argument uses neither a choice of representatives nor
a splitting assumption.

\begin{definition}[Regular category]\label{f:regular}
A category is \emph{regular} if it has finite limits, every map factors
as a regular epimorphism followed by a monomorphism, and pullbacks of
regular epimorphisms are regular epimorphisms. The monomorphism in
\[
X\twoheadrightarrow \im(f)\hookrightarrow Y
\]
is called the \emph{image} of $f:X\to Y$.
\end{definition}

\begin{lemma}[Image calculus]\label{f:images}
In a regular category, images are unique as subobjects, regular
epimorphisms compose, and pulling back an image factorization gives an
image factorization.
\end{lemma}
\begin{proof}
Given two regular-epimorphism--monomorphism factorizations of a map,
Lemma~\ref{f:descent} gives a map between their middle objects in each
direction. Its uniqueness says the two composites are identities.
This proves uniqueness of images.

The lifting property in Lemma~\ref{f:descent} is preserved under
composition: first lift against the first regular epimorphism and then
against the second. Factor their composite as $mp$, where $p$ is a
regular epimorphism and $m$ is monic. The lifting property of the
composite, applied against $m$, gives a right inverse of $m$. A monic
map with a right inverse is an isomorphism. The composite is therefore
a regular epimorphism.

Pull back the mono part of a factorization, and then pull back its
regular epi part. The resulting square is a pullback of the composite,
by the defining universal properties. Pullback stability makes the
first part a regular epimorphism; pullbacks of monomorphisms are monic
because maps into the pullback are determined by their two projections.
This is the required image factorization.
\end{proof}

Images express what survives a map. Pullback stability says that this
description remains valid after imposing a compatible condition at the
target. It is the reason local witnesses can safely be used below.

\subsection{A precise license to reason with witnesses}

An arbitrary category may have too few maps from its terminal object to
detect equality. In particular, a group has just one group homomorphism
from the trivial group. Thus an argument about ordinary global points
would miss nearly everything. We instead use \emph{generalized elements}:
an input $x$ of $X$ over a test object $T$ means a map $x:T\to X$.

For a subobject $M\hookrightarrow X$, saying that $x$ belongs to $M$
means that $x$ factors through that monomorphism. A \emph{cover} below
means a regular epimorphism, not an assumed topological covering.

\begin{lemma}[Local witness calculus]\label{f:local-calculus}
Let $f:A\to X$ be a map in a regular category. A generalized element
$x:T\to X$ belongs to $\im(f)$ if and only if there is a cover
$e:T'\twoheadrightarrow T$ and a map $a:T'\to A$ such that $fa=xe$.
Finite lists of such witnesses can be chosen over a common cover.
They can be pulled back along any map of test objects.
\end{lemma}
\begin{proof}
Factor $f=mp$ with $p:A\twoheadrightarrow I$ regular epic and $m:I\hookrightarrow X$
monic. If $x=mi$, pull back $p$ along $i:T\to I$. The pullback cover
has a map $a$ with $fa=xe$. Conversely, such $e,a$ give a square
against $m$; Lemma~\ref{f:descent} makes $x$ factor through $m$.

Pulling back a cover gives a cover. To combine two choices over $T$,
take their fiber product over $T$; either projection is a cover, and
its composite to $T$ is a cover by Lemma~\ref{f:images}. Repeat for
finitely many choices. Nested choices are handled by composing covers.
\end{proof}

This is the precise meaning of ``there exists'' in our relation
calculations. A witness may exist only after a cover of the test object.
Every proof using witnesses below permits that replacement. Equalities
of subobjects are checked by inclusions in both directions: to check
$M\leq N$, one may test the generalized element $M\hookrightarrow X$
itself. The method therefore proves actual categorical statements,
not merely assertions about points.

\subsection{Relations compose by matching, then forgetting}

\begin{definition}[Internal relation]
A relation $R:X\rightsquigarrow Y$ is a subobject of $X\times Y$.
Choose a representative with projections
\[
\begin{tikzcd}[column sep=large]
& M \arrow[dl,"a"'] \arrow[dr,"b"] &\\
X && Y.
\end{tikzcd}
\]
The pair $(a,b)$ is jointly monic. We may write $x\,R\,y$ for membership
in this subobject. Its opposite $R^{\mathrm{op}}:Y\rightsquigarrow X$
interchanges the projections. The identity relation $\Delta_X$ is the
diagonal $X\hookrightarrow X\times X$.
\end{definition}

If $S:Y\rightsquigarrow Z$ has carrier $N$, form $M\times_YN$ and take
its image in $X\times Z$. This is the composite $SR$. Thus
\[
x\,(SR)\,z
\quad\Longleftrightarrow\quad
\text{locally there is }y\text{ with }x\,R\,y\text{ and }y\,S\,z.
\]
\emph{Our convention is always that $SR$ means first $R$, then $S$.}
The pullback retains the matching intermediate value; the image forgets
which matching value was used.

\begin{proposition}[The calculus of relations]\label{f:relation-calculus}
In a regular category, relation composition is associative and has
the diagonal relations as identities. It preserves inclusions, and
\[
(SR)^{\mathrm{op}}=R^{\mathrm{op}}S^{\mathrm{op}}.
\]
\end{proposition}
\begin{proof}
For $R:X\rightsquigarrow Y$, $S:Y\rightsquigarrow Z$, and
$T:Z\rightsquigarrow W$, membership in either $(TS)R$ or $T(SR)$
means that, over a cover, there are $y,z$ with
$x\,R\,y$, $y\,S\,z$, $z\,T\,w$. Lemma~\ref{f:local-calculus}
combines any nested witnesses into one cover and also allows the same
triple witness to be grouped in either order. This proves equality of
the two image subobjects. Equivalently, both are the image of
$M_R\times_YM_S\times_ZM_T$ in $X\times W$.

A diagonal intermediate relation forces the two matched coordinates
to agree and introduces no further condition; hence the unit laws.
An inclusion of relations sends every witness to a witness, proving
monotonicity. Reversing a witness exchanges its endpoints and reverses
the order of the two conditions, proving the last identity.
\end{proof}

This associativity has the same shape as elementary reassociation, but
it applies to correspondences that may not be functions. Its proof
identifies one object of complete witness data and shows that the order
of forgetting intermediate information does not matter.

A map $f:X\to Y$ is identified with its graph relation, represented by
$(\id_X,f):X\hookrightarrow X\times Y$. The opposite graph is generally
not a map: it remembers all possible preimages.

\begin{lemma}[Graphs, kernels, and supports]\label{f:graph-kernel}
For a map $f:X\to Y$,
\[
f^{\mathrm{op}}f=\operatorname{Eq}(f),
\qquad
ff^{\mathrm{op}}=\Delta_{\im(f)}
\quad\text{as a relation on }Y.
\]
Here $\operatorname{Eq}(f)=X\times_YX$ is the kernel pair, and the second
diagonal is supported on the image subobject of $Y$. Consequently
\[
ff^{\mathrm{op}}f=f,
\qquad ff^{\mathrm{op}}=\Delta_Y
\text{ if }f\text{ is regular epic}.
\]
Every relation with projections $a:M\to X$, $b:M\to Y$ satisfies
\[
R=ba^{\mathrm{op}}.
\]
\end{lemma}
\begin{proof}
The first composite relates $x,x'$ exactly when $f(x)=f(x')$, which is
the defining pullback. The second relates $y,y'$ exactly when, locally,
$y=f(x)=y'$ for some $x$; this is the image diagonal by
Lemma~\ref{f:local-calculus}. These descriptions give the two
consequences. Finally $ba^{\mathrm{op}}$ is the image of $(a,b)$ in
$X\times Y$; that pair already represents the subobject $R$.
\end{proof}

\begin{lemma}[When a relation is a map]\label{f:functional}
A relation $R:X\rightsquigarrow Y$ is the graph of a unique map
$f:X\to Y$ if
\[
\Delta_X\leq R^{\mathrm{op}}R,
\qquad RR^{\mathrm{op}}\leq\Delta_Y.
\]
The first condition says that every input has a locally available
output; the second says that the output is unique.
\end{lemma}
\begin{proof}
Write $R=ba^{\mathrm{op}}$ with carrier $M$. The first condition says
that the image of $a:M\to X$ is all of $X$, so $a$ is regular epic.
If $u,v:T\to M$ have $au=av$, then $bu$ and $bv$ are two outputs for
the same input. The second condition gives $bu=bv$. Joint monicity
of $(a,b)$ gives $u=v$, so $a$ is also monic and hence an isomorphism.
Now $f=ba^{-1}$ has graph $R$. Its value is forced by that graph,
which proves uniqueness.
\end{proof}

\subsection{Exactness makes an identification into an object}

A relation $E:X\rightsquigarrow X$ is \emph{reflexive} if
$\Delta_X\leq E$, \emph{symmetric} if $E^{\mathrm{op}}=E$, and
\emph{transitive} if $EE\leq E$. An \emph{internal equivalence relation}
has all three properties. Reflexivity also gives $E\leq EE$, so an
equivalence relation satisfies $EE=E$.

Kernel pairs are equivalence relations. For example, transitivity says
that $f(x)=f(y)$ and $f(y)=f(z)$ imply $f(x)=f(z)$; this is equally an
equality of maps over any test object.

\begin{definition}[Barr exactness]\label{f:exact}
A regular category is \emph{Barr-exact} if every internal equivalence
relation is the kernel pair of a map. Equivalently, for every such
$E$ on $X$ there is a regular epimorphism $q:X\twoheadrightarrow Q$
with $E=q^{\mathrm{op}}q$. We call $Q$ the quotient $X/E$.
\end{definition}

For the equivalence of these formulations, factor a map exhibiting the
kernel pair into its image. The monomorphism part does not change its
kernel pair, and the other part is regular epic. Lemma~\ref{f:descent}
then gives the quotient's full universal property: maps from $X$ that
identify the pairs in $E$ descend uniquely to $Q$. Two such quotients
are uniquely isomorphic over $X$, by applying this property twice.
Exactness concerns these equivalence-relation quotients; it does not
assert that the category is abelian or that its maps can be added.

A category is \emph{pointed} if it has an object $0$ that is both initial
and terminal. There is then a zero map between every two objects, namely
the composite through $0$. A \emph{kernel} of $f:X\to Y$ is the
pullback of $0\to Y$ along $f$:
\[
\begin{tikzcd}
\Ker(f) \arrow[r] \arrow[d,hook] & 0 \arrow[d]\\
X \arrow[r,"f"'] & Y.
\end{tikzcd}
\]
The left map is monic since $0\to Y$ is monic. In a pointed category
it is a distinguished subobject, but it is only one class of the whole
kernel pair.

If $E=q^{\mathrm{op}}q$, its zero class is the relation $E0_X$,
viewed as a subobject of $X$, where $0_X:0\to X$. Since $q0_X=0_Q$,
\[
E0_X=q^{\mathrm{op}}0_Q=\Ker(q).
\]
It is saturated: $E(E0_X)=E0_X$. We will retain both the full
equivalence relation and its zero class. In groups the zero class
determines the normal quotient, but we will not build that additional
fact into the general framework.

The roles of the assumptions are now visible. Finite limits impose
compatible equations. Images forget witnesses. Regularity makes that
forgetting compatible with substitution. Exactness realizes consistent
identifications as quotient objects. Pointedness selects a zero class.
The next condition will make different identifications compatible with
one another.
